\section{Proofs for Section~\ref{sec:projection}}\label{app:projection}

\subsection{Proof of Theorem~\ref{thm:projection}}
Fix one allowed matching $M$ with $|M|=m$, fleet $N=N_m$, requirement
$R=\sum_{(i,j)\in M}(a_i+b_j-B)_+$ and $T=S_0-NB$. Let $x_v$ be a bus's net early
battery gain. A paired bus serving $A_i$ and $B_j$ holds $B-a_i$ after its
morning service; it cannot exceed capacity, so $x_v\le a_i$, and it must hold at
least $b_j$ before the afternoon service, so $x_v\ge a_i+b_j-B$. An A-only bus
has $x_v\in[a_i-B,a_i]$, and a B-only bus, full before the early window, has
$x_v\in[b_j-B,0]$. These boxes are nonempty because $a_i,b_j\le B$. Their lower
endpoints sum to $T$, their positive lower endpoints to $R$, and their upper
endpoints to $P$; every upper endpoint is nonnegative.

\emph{Necessity.} Whatever the order of actions, $x_v=\alpha C_v-X_v/\delta$ for
the bus's early purchases $C_v$ and exports $X_v$, so the total $Q$ lies in
$[T,P]$. A positive lower endpoint needs at least its value divided by $\alpha$
in purchases, hence $C\ge R/\alpha$. Exports are nonnegative, the shared
connector carries gross energy $C+X\le K$, and full terminal replenishment
needs the terminal purchase $L=(S_0-Q)/\alpha\le H$. Terminal purchases are
nonnegative because $Q\le P<S_0$.

\emph{Sufficiency.} Fix totals $(C,X)$ satisfying these inequalities with $R$ in
place of $r_m$. For any $Q\in[T,P]$ the boxes admit gains whose positive parts
sum to $A=\max\{R,Q\}$. If $Q\ge R$, start from the positive lower endpoints and
zero elsewhere, and distribute $Q-R$ through the remaining headroom. If $Q<R$,
keep the mandatory positive gains and distribute the negative total $Q-R$
through the negative lower intervals, whose combined lower endpoint is $T-R$.
One-direction exchanges then realize $C_0=A/\alpha$ and $X_0=\delta(A-Q)$. Since
$X\ge0$ implies $\alpha C\ge Q$, and $\alpha C\ge R$, the residual $z=C-C_0$ is
nonnegative and $X-X_0=\alpha\delta z$.

Realize these residuals by battery-neutral cycles on one bus before its
one-direction exchange. At early state of charge $s\in[0,B]$, use
charge-then-discharge cycles of depth at most $B-s$ if $s\le B/2$, and
discharge-then-charge cycles of depth at most $s$ otherwise. The permitted depth
is at least $B/2>0$. Split the residual battery throughput $\alpha z$ into
finitely many such chunks. Each returns to $s$, stays in $[0,B]$, and uses grid
purchases depth$/\alpha$ and exports $\delta\cdot$depth, so together they
realize exactly $(z,\alpha\delta z)$. Execute the cycles and then all
one-direction exchanges serially: their gross energy is $C+X\le K$, so they fit
the common early window at the connector's power. The endpoint boxes guarantee
the afternoon service, and nonnegative service draw preserves every
intermediate state-of-charge bound. Refill serially at the terminal connector
using total $L\le H$. If $K=0$ all early actions vanish; if $H=0$, then $P<S_0$
makes every polygon empty, consistent with infeasibility.

\emph{Replacement.} At fixed $m$, the quantities $S_0,P,N_m,T$ and the fleet cost
do not depend on the matching; only $R$ does, and only through $C\ge R/\alpha$.
Hence the polygon of a matching attaining $r_m$ contains the polygon of every
other matching of size $m$, and every point of it has just been realized.
The union over matchings of size $m$ therefore equals the polygon $W_m$, and the
union over $m$ gives $\mathcal P$. Each $W_m$ is defined by five half-planes, so
it has at most five vertices, and $\conv\mathcal P$ is the convex hull of their
images. With exports prohibited, the construction above uses no exports when
$X=0$: then $\alpha C=Q\ge R$, so $A=Q$, $X_0=0$ and $z=0$. \qed

\paragraph{Idle buses.} Identical paid idle buses under the total fleet cap
$n_A+n_B$ do not change $\mathcal P$. An idle bus has the box $[-B,0]$. A plan
with $m$ pairs and $q$ idle buses has $q\le m$ by the cap; deleting $q$ pairs
and serving their services separately preserves $N,S_0,P,T$ and the fleet cost
while lowering $R$, and reproduces the same $(C,X)$.

\subsection{The lossless cases}
\emph{Charge-only.} With $\alpha=\delta=1$ and $X=0$, $W_m$ becomes
$\{C:\ C\ge r_m,\ C\le K,\ S_0-N_mB\le C\le P,\ C\ge S_0-H\}$. Since
$S_0-N_mB$ is the sum of all lower endpoints and $r_m$ the sum of their positive
parts, $S_0-N_mB\le r_m$, and the layer is $\{E:\max\{\mathrm{Lo},r_m\}\le E\le U\}$ with
$L=S_0-E$, which is \eqref{eq:lossless}. Feasibility is $\mathrm{Lo}\le U$ because the empty
matching has $r_0=0$.

\emph{Bidirectional.} With $\alpha=\delta=1$ and exports allowed, $E=C-X$ and
$L=S_0-E$. A point $E$ is attainable in layer $m$ if and only if
$\max\{S_0-N_mB,S_0-H\}\le E\le P$ and some $C\ge\max\{r_m,E\}$ satisfies
$2C-E\le K$, that is, $E\le K$ and $E\ge2r_m-K$. The layer is therefore
$[t_m,U]$ with $t_m=\max\{S_0-H,\,S_0-N_mB,\,2r_m-K\}$ and $U=\min\{P,K\}$.

\subsection{Monotonicity, convexity and the frontier}
Deleting an edge from a minimum matching of size $m$ shows $r_{m-1}\le r_m$.
For discrete convexity, consider the unit-capacity source--A--B--sink network on
the allowed edges with costs $(a_i+b_j-B)_+$. Its minimum-cost flow LP is
integral at integer flow values $m=0,\dots,\nu$, and its value is convex in the
required flow, so $2r_m\le r_{m-1}+r_{m+1}$. This uses bipartiteness, not
completeness of $\Gamma$.

All thresholds at or below $\mathrm{Lo}$ are dominated in cost by $m_0$. The first retained
increment is $r_{m_0+1}-\mathrm{Lo}\le r_{m_0+1}-r_{m_0}$, and later increments are the
original nondecreasing matching increments, all positive after $m_0$. The slopes
$-f/\delta_m$ are therefore nondecreasing, followed by the flat segment.
Endpoints are physical, interpolants are hull mixtures, and every physical
branch lies on or above this convex graph, so the graph is exactly the lower
hull boundary. Consecutive breakpoints differ by one cardinality.

\subsection{Proof of Corollary~\ref{cor:support}}
The convex frontier plus strictly convex continuous $G$ has a unique minimizer
$E^*$. At $\mathrm{Lo}$, at a breakpoint or on the flat segment, a physical plan attains
the frontier, so $D=\CH$. Strictly between consecutive breakpoints the frontier's
operating cost is strictly below that of every physical branch because $f>0$, so no
physical plan attains the hull optimum; compactness of the finite union of
physical segments then gives $D>\CH$.

\subsection{Sorted matching for the complete graph}
Replacing a selected service by an unused smaller draw while keeping its partner
cannot increase the nondecreasing edge cost, so some optimal size-$m$ matching
uses the $m$ smallest draws in each batch. With $h(t)=(t-B)_+$ and $a\le a'$,
$b\le b'$, convexity gives $h(a+b)+h(a'+b')\ge h(a+b')+h(a'+b)$. Swapping any
two partners ordered in the same direction therefore cannot raise the cost, and
repeated exchanges yield opposite order. For a non-complete graph this shortcut
fails: with A draws $(5,15)$, B draws $(14,15)$ and $B=20$, allowing only the
index-matched edges gives $r_2=10$ instead of $9$.

\subsection{The worked example}
The requirement matrix is $\bigl(\begin{smallmatrix}0&0\\9&10\end{smallmatrix}\bigr)$,
so $r_0=r_1=0$, $r_2=9$, $S_0=49$, $\mathrm{Lo}=0$ and $U=15$. With
$G(E)=2401/20-9E/10+E^2/10$, the physical minima of layers $m=0,1$ are at $E=9/2$,
with values $5841/40$ and $5561/40$; layer $m=2$ is minimized at its endpoint
$E=9$ with value $2681/20$. On the sloping hull face the objective is
$2821/20-151E/90+E^2/10$, minimized at $E^*=151/18$ with value $434201/3240$, so
$\Delta=121/3240$. The hull point mixes the three-bus plan at $E=0$ and the
two-bus plan at $E=9$ with weights $11/162$ and $151/162$. At the physical
optimum's marginal prices $(49/10,4)$ each layer's bill is minimized at its lowest
endpoint; bill plus vehicle cost is $224$, $217$ and $2181/10$, so the regret is
$11/10$. Both mixed plans have explicit serial witnesses: at $E=0$ pair the 5-draw
A with the 14-draw B and refill $19,15,15$ on the terminal connector; at $E=9$ the
15-draw A bus charges $9$ at midday in the crossed pairing and both buses refill
$20$.

\section{Proof of Proposition~\ref{prop:support}}\label{app:support}
For any hull point $(e,c)$, linear minimization gives $c+p^\top e\ge V(p)$. If
$p\in\partial F(e(x))$, then $F(e)\ge F(e(x))+p^\top(e-e(x))$. Adding and
minimizing over the hull gives $\CH\ge V(p)+F(e(x))-p^\top e(x)=c(x)+F(e(x))-r(x;p)$,
which is the regret bound; if $r(x;p)=0$, then $D=\CH$.

Conversely let $D=\CH$. Compactness gives a physical minimizer $x$. The
relative-interior condition gives strong Fenchel duality with an attained
maximizer, $\CH=V(p^*)-F^*(p^*)$. Then
\[
 \underbrace{c(x)+(p^*)^\top e(x)-V(p^*)}_{\ge0}
 +\underbrace{F(e(x))-(p^*)^\top e(x)+F^*(p^*)}_{\ge0}=D-\CH=0,
\]
so both terms vanish: Fenchel equality gives $p^*\in\partial F(e(x))$, and the
first term gives $r(x;p^*)=0$. \qed

\section{Proofs for Section~\ref{sec:prices}}\label{app:prices}

Throughout, $\bar c$ is the lower hull boundary on $[\mathrm{Lo},U]$, extended by
$+\infty$ outside, and $H(\pi)=\arg\min_E\{\bar c(E)+\pi E\}$ is the hull response.
The finite candidate set $\mathcal V$ consists of all layer endpoints $(\ell_m,fN_m)$
and $(U,fN_m)$. Since a linear function attains its minimum over $\conv\mathcal P$
at a point of $\mathcal V$, every $E\in R(\pi)$ satisfies $\bar c(E)=fN_m$ for its
layer and $E\in H(\pi)$. The layer endpoints that lie on $\bar c$ are exactly the
breakpoints $\mathrm{Lo}$, $r_m$ ($m_0<m\le m_{\max}$) and $U$. The hull optimum satisfies
$-\pi^*\in\partial\bar c(E^*)$ because $G$ is differentiable.

\subsection{Proof of Proposition~\ref{prop:anatomy}}
(i) By \cref{cor:support}, $\Delta>0$ places $E^*$ strictly between consecutive
breakpoints $u<u'$, where $\bar c$ is affine with slope $-f/(u'-u)$ because
consecutive breakpoints differ by one cardinality. Optimality at an interior point
gives $g(E^*)=f/(u'-u)\ge f/\delta_{\max}$, and $\max g\ge g(E^*)$.

(ii) On $[u,u']$ let $\phi=\bar c+G$. Then $\phi'(E)=-f/(u'-u)+g(E)$ is
$\kappa$-Lipschitz and vanishes at $E^*$, so $\phi(E)-\phi(E^*)\le\kappa(E-E^*)^2/2$.
Both $u$ and $u'$ are physical with cost $\bar c$, so
$D\le\min\{\phi(u),\phi(u')\}$ while $\CH=\phi(E^*)$. One of $u,u'$ is within
$(u'-u)/2$ of $E^*$, which gives $\Delta\le\kappa(u'-u)^2/8$. \qed

\subsection{Proof of Lemma~\ref{lem:sign}}
Let $E\in R(\pi)\subseteq H(\pi)$, so $-\pi\in\partial\bar c(E)$, and recall
$-\pi^*\in\partial\bar c(E^*)$. Monotonicity of the subdifferential of the convex
function $\bar c$ gives $(-\pi+\pi^*)(E-E^*)\ge0$, that is, $(\pi-\pi^*)(E-E^*)\le0$.
Since $g$ is increasing and $g(E^*)=\pi^*$, $(\pi-\pi^*)(g(E)-\pi^*)\le0$. For the
second claim write
$\pi_{k+1}-\pi^*=(1-a)(\pi_k-\pi^*)+a(g(E_k)-\pi^*)$. The two terms have
opposite (weak) signs, and for real $x,y$ with $xy\le0$, $|x+y|\le\max\{|x|,|y|\}$.
Since $E_k\in[\mathrm{Lo},U]$, $|g(E_k)-\pi^*|\le M$. \qed

\subsection{Proof of Theorem~\ref{thm:iteration}}
We use three facts. (F1) The same monotonicity argument shows that $\pi<\pi'$,
$E\in H(\pi)$ and $E'\in H(\pi')$ imply $E\ge E'$. (F2) If $E^*$ lies in the relative
interior of a maximal face $[w,w']$, then $\partial\bar c(E^*)=\{-\pi^*\}$, the face
has slope $-\pi^*$, $[w,w']\subseteq H(\pi^*)$, and $E^*\in H(\pi)$ only for
$\pi=\pi^*$. By (F1), $H(\pi)\subseteq[\mathrm{Lo},w]$ for $\pi>\pi^*$ and
$H(\pi)\subseteq[w',U]$ for $\pi<\pi^*$. (F3) A fixed point $\hat\pi=g(\hat E)$ with
$\hat E\in R(\hat\pi)$ has $-g(\hat E)\in\partial\bar c(\hat E)$, so $\hat E$
minimizes the convex function $\bar c+G$; by uniqueness $\hat E=E^*$ and
$\hat\pi=\pi^*$.

\emph{(a)} By (F3), a fixed point requires $E^*\in R(\pi^*)$ to be a selectable
layer endpoint on $\bar c$, that is, a breakpoint. Conversely, if $E^*$ is a
breakpoint then $E^*\in H(\pi^*)$ is a candidate attaining the minimum, so
$E^*\in R(\pi^*)$ and a selection choosing it makes $\pi^*$ a fixed point. If
$\Delta>0$, $E^*$ is not physical and hence not a breakpoint; if $E^*$ lies inside
the flat face $(\ell_{m_{\max}},U)$ it is physical but not a breakpoint.

\emph{(b)} Let $\sigma$ be the selection, a function of $\pi$ with values in
$R(\pi)$; by (F1) it is nonincreasing, so $h=g\circ\sigma$ is nonincreasing and
$h\circ h$ is nondecreasing. For $k\ge1$, $\pi_k$ lies in the finite set
$g(\{E:(E,c)\in\mathcal V\})$. The odd and even subsequences are orbits of the
nondecreasing map $h\circ h$, hence monotone, hence eventually constant in a finite
set. Thus $\pi_{k+2}=\pi_k$ eventually. In the face case with $\pi_0\neq\pi^*$, we
show by induction that $\pi_k\ne\pi^*$: if $\pi_k\neq\pi^*$ then by (F2)
$E_k\ne E^*$, so $\pi_{k+1}=g(E_k)\ne\pi^*$. By (F3) there is then no fixed
point along the orbit, and the period is two. If $\pi_a>\pi^*$ is one state of the
cycle, (F2) gives $E_a\le w<E^*$, hence $\pi_b=g(E_a)<\pi^*$ and $E_b\ge w'$.

\emph{(c)} By \cref{lem:sign}, $d_k=|\pi_k-\pi^*|$ satisfies
$d_{k+1}\le\max\{(1-a)d_k,aM\}$, so $d_k\le\max\{(1-a)^kd_0,aM\}$ and the
interval $d\le aM$ is reached in finitely many steps and is invariant.

\emph{Vertex case.} If $\pi^*$ is interior to the normal cone
$\{\pi:E^*\in H(\pi)\}$ and $aM$ is less than its distance to the cone's boundary,
then inside the invariant interval $H(\pi)=\{E^*\}$, so $E_k=E^*$ and
$\pi_{k+1}-\pi^*=(1-a)(\pi_k-\pi^*)$.

\emph{Face case: non-convergence.} Suppose $\pi_k\to\bar\pi$. Then
$a(g(E_k)-\pi_k)=\pi_{k+1}-\pi_k\to0$, so $g(E_k)\to\bar\pi$. Since $E_k$ ranges over a
finite set and $g$ is injective, $E_k=\bar E$ for all large $k$, with
$\bar\pi=g(\bar E)$. The set $\{\pi:\bar E\in R(\pi)\}$ is closed, being defined by
finitely many weak linear inequalities, so $\bar E\in R(\bar\pi)$, and (F3) gives
$\bar E=E^*$ and $\bar\pi=\pi^*$. If $E^*$ is not a breakpoint it is never selected,
which is a contradiction. If it is a breakpoint, $E^*\in R(\pi_k)$ forces
$\pi_k=\pi^*$ by (F2), so the orbit hits $\pi^*$ exactly. For a fixed sequence of
selected loads, $\pi_k$ is affine in $\pi_0$ with slope $(1-a)^k\neq0$, so each
finite selection sequence admits at most one such $\pi_0$, and there are countably
many sequences.

\emph{Face case: trapping.} Write $[\pi^*,\sigma]$ for the normal cone of $w$ and
$[\sigma',\pi^*]$ for that of $w'$ (with $\sigma=+\infty$ if $w=\mathrm{Lo}$ and
$\sigma'=-\infty$ if $w'=U$), and assume $aM<\min\{\sigma-\pi^*,\pi^*-\sigma'\}$.
Inside the invariant interval, $\pi>\pi^*$ forces $E=w$ and $\pi<\pi^*$ forces
$E=w'$; at $\pi=\pi^*$ the selection is a breakpoint in $[w,w']$. With
$d_-=\pi^*-g(w)>0$ and $d_+=g(w')-\pi^*>0$, an update from $\pi>\pi^*$ gives
$\pi'=(1-a)\pi+ag(w)\in[\pi^*-ad_-,\pi)$, and symmetrically from below; from
$\pi^*$ the update stays in $[\pi^*-ad_-,\pi^*+ad_+]$. This band is therefore
invariant. Above $\pi^*$, prices contract geometrically toward $g(w)<\pi^*$ and cross
$\pi^*$ in finitely many steps; below, they contract toward $g(w')>\pi^*$. Hence
the band is entered and both $w$ and $w'$ recur infinitely often.

\emph{(d)} With $d_{k+1}\le\max\{(1-a_k)d_k,a_kM\}$, fix $\varepsilon>0$ and $K$ with
$a_kM<\varepsilon$ for $k\ge K$. After $K$, the interval $d\le\varepsilon$ is invariant,
and outside it $d$ contracts by the factors $1-a_k$, whose product diverges to
zero because $\sum a_k=\infty$. Thus $\limsup d_k\le\varepsilon$ for every
$\varepsilon$. Summing the update gives $\sum_{k<K}a_k(g(E_k)-\pi_k)=\pi_K-\pi_0$,
which is bounded; dividing by $A_K=\sum_{k<K}a_k\to\infty$ and using
$\pi_k\to\pi^*$ (Stolz--Ces\`aro) gives $\sum_{k<K}a_kg(E_k)/A_K\to\pi^*$. If $g$
is affine, this average equals $g\bigl(\sum_{k<K}a_kE_k/A_K\bigr)$, and invertibility
of $g$ gives the load statement. \qed

\paragraph{Numerical verification.} The repository contains an exact
rational-arithmetic implementation of the projection, hull, gap and iteration
(\texttt{code/twobatch.py}) and randomized checks of every statement in this
section on several thousand instances: the period and bracketing claims in (b);
the absence of fixed points; trapping, band and recurrence in (c) at half the
stated step threshold; convergence in (d); and both parts of
\cref{prop:anatomy}. Instances are drawn so that the hull optimum falls on
sloping faces, flat faces and vertices.

\section{Proofs for Section~\ref{sec:boundaries}}\label{app:boundaries}

\subsection{Common transition energy}
Suppose every matched pair performs a compulsory task of energy $d\ge0$, with
nonnegative draw, during a noncharging interval after the early window and
before the B batch. A paired bus then needs $x_v\ge a_i+b_j+d-B$ to finish its
afternoon service, cannot exceed $x_v\le a_i$, and needs terminal refill
$a_i+b_j+d-x_v$. The box is nonempty exactly when $b_j+d\le B$, so edges
violating this are removed. Lower endpoints sum to $S_m-N_mB$ with
$S_m=S_0+md$, positive lower endpoints to the matching's requirement with $d$
included, and the upper endpoints still to $P$. Because $S_m$ depends only on $m$,
the proof of \cref{thm:projection} applies verbatim with $S_0$ replaced by $S_m$.

\subsection{Proof of Proposition~\ref{prop:completegraph}}
\emph{A disjoint-block construction.} Let $w_i=2^i$ for $i=0,\dots,k-1$,
$W=\sum_iw_i$, $B=20\max_iw_i$. Block $i$ has two A and two B services with draws
$(B/2+3w_i,\,w_i)$ in each batch. Use windows A $[0,1]$, early charging $[1,2]$,
noncharging transition tasks $[2,3]$, B $[3,4]$ and terminal charging $[4,5]$; one
early connector of budget $K=kB/2+4W$ and one terminal connector. Within a block,
the identity edges carry transition energy $0$ and the cross edges $6w_i$. The fleet
cap is $2k$, so exactly $2k$ buses are used and every bus serves one service from
each batch. Identity requires early purchase $6w_i$ and consumes $B+8w_i$; cross
requires none and consumes $B+20w_i$. For a set $I$ of identity blocks, with
$R=6\sum_{i\in I}w_i$ and $S_{\max}=kB+20W$, total consumption is $S_{\max}-2R$, the
per-bus early intervals sum to $[R,K]$, and the complete branch is
$\{(E,\,S_{\max}-2R-E,\,2kf):R\le E\le K\}$, each point realized by serial charging.

\emph{Complete graph.} Now admit every A--B edge. Keep the within-block transition
energies, and on each edge between different blocks set $d_{ij}=B-b_j$; raise the
terminal budget to $H=2kB$. For any edge put $e_{ij}=a_i+b_j+d_{ij}$ and
$r_{ij}=(e_{ij}-B)_+$. Its early purchase interval is $[r_{ij},a_i]$ whenever
$b_j+d_{ij}\le B$, which holds for within-block edges because $B\ge20\max w_i$ and
with equality between blocks. In the latter case $e_{ij}=a_i+B$ and $r_{ij}=a_i$: the
bus refills completely at midday. For a perfect matching $M$ these intervals sum
to $[R_M,K]$ with $R_M=\sum_Mr_{ij}$, total consumption is $S_M=\sum_Me_{ij}$, and the
branch is $L=S_M-E$, $R_M\le E\le K$. Each terminal purchase is at most
$e_{ij}-r_{ij}\le B$, so $H=2kB$ accommodates it; $E=K$ is feasible for every
perfect matching.

Give each high service (draw $B/2+3w_i$) in block $i$ the potential $B/2+9w_i$ and
each low service the potential $w_i$, as row potentials $u_i$ and column potentials
$v_j$; their total is $S_{\max}$. Every within-block edge satisfies
$e_{ij}+2r_{ij}=u_i+v_j$. A between-block edge has $e_{ij}+2r_{ij}=B+3a_i$, with slack
over $u_i+v_j$ equal to $3B/2-9w_j$, $2B-w_j$, $B/2+2w_i-9w_j$ or $B+2w_i-w_j$ for the
four high/low combinations. Each is at least $\gamma=w_{\max}+2w_{\min}\ge3$, the
low-to-high case being limiting because $B/2\ge10w_{\max}$. Summing this
assignment-potential certificate over $M$ gives $S_M+2R_M\ge S_{\max}$, with surplus at
least $\gamma$ if any between-block edge is used. Hence every physical point, and
every mixture, satisfies $3E+L=S_M+2E\ge S_M+2R_M\ge S_{\max}$.

With $F(E,L)=3E+L+(E-3)^2/(2K)$, total cost is at least $2kf+S_{\max}$. The equal
mixture of the within-block corners $R=0$ and $R=6$ attains it at load
$(3,S_{\max}-9)$, so $\CH=S_{\max}+2kf$. On a within-block branch the cost is
$2kf+S_{\max}+2(E-R)+(E-3)^2/(2K)$, increasing in $E$ (derivative at least
$2-3/K>0$ since $K\ge14$), hence minimized at $E=R$; the values $R=6t$ nearest to 3
are $0$ and $6$, giving excess $9/(2K)$. Every between-block schedule has excess at
least $\gamma\ge3>9/(2K)$. Therefore $D=\CH+9/(2K)$. \qed

\subsection{Proof of Proposition~\ref{prop:threebatch}}
Reduce from numerical three-dimensional matching, strongly NP-complete
\citep[SP16]{garey1979}: three multisets of $n$ positive integers must be
partitioned into triples, one from each, each summing to $B$. We may assume total
$nB$ and every number at most $B$, mapping violations to a fixed NO instance. Make
the multisets services in $[0,1)$, $[1,2)$ and $[2,3)$ with those numbers as
battery draws, with charging unavailable before time 3. Overlap and the fleet cap
force exactly $n$ buses, each serving one service per batch; a duty $(a,b,c)$ is
feasible exactly when $a+b+c\le B$, and the total $nB$ forces equality. Conversely, a
matching yields duties ending empty, which refill serially for $B$ time units each
during $[3,3+nB]$. The construction is polynomial and preserves strong hardness;
membership in NP is immediate. \qed

\section{The certified public timetable}\label{app:public}
The case uses the Hildenbrand timetable of \citet{sistig2025dataset}: 37 mandatory
services, the complete directed $5\times5$ deadhead matrix, and depot 15 as the sole
charging and parking location. Service energy is $1.2$\,kWh/km plus 16\,kWh per
service hour; deadheads use $0.96$\,kWh/km plus 16\,kWh per hour; off-depot waiting
uses 16\,kWh/h. Each used bus has 400\,kWh usable energy and 360\,kW acceptance; one
360\,kW connector is shared; every bus starts full and returns full by 30:00. The
used-bus cost is 100 and the supply cost is
$F(e)=\sum_{t=1}^{30}(0.2e_t+e_t^2/1800)$ over hourly intervals, with the exact binary
values of these coefficients used throughout.

\emph{Upper bound on $\CH$.} An incumbent physical plan and a second admitted
physical response, which has a higher planning cost but a lower bill at the
incumbent's marginal prices, are mixed. Exact quadratic evaluation and
one-dimensional optimality conditions give $\CH\le441.142238$.

\emph{Lower bound on $D$.} A minimum-fleet argument excludes one bus. For three or
more buses, vehicle cost plus the linear mandatory-service energy term already
exceeds $477.82$. The two-bus domain is covered by a binary branching tree over
movement variables. At each leaf, exact tangent planes, signed row multipliers and
finite-box residual compensation give a valid rational lower bound, and infeasible
leaves carry rational infeasibility certificates. An independent reconstruction
checks coverage and every leaf, giving $D\ge441.532912$. A feasible schedule gives
$D\le441.532913$.

\emph{Scope.} The first certificate concerns the declared movement graph (pull-out,
direct, via-depot and pull-in modes with fixed leg times). A second, separately
checked certificate uses an exactly additive energy law and, under directed
triangle inequalities that the source tables satisfy, extends the result to
arbitrary finite walks over the location pairs with continuous waiting. The
displayed bounds coincide, but neither certificate is transferred to the other model. The
gap holds for every vehicle cost $f\ge100$, since both mixture components and the
two-bus lower bound shift equally. Standalone checkers using only the Python standard
library re-verify each certificate in under a minute (\cref{app:repro}).

\section{Reproducibility}\label{app:repro}
All code and derived data are in the repository at \texttt{paper/draft1/}:
\begingroup\sloppy
\begin{itemize}[itemsep=1pt]
\item \texttt{code/twobatch.py}: exact projection, hull, gap and price iteration;
\texttt{code/test\_twobatch.py}: the randomized checks of \cref{sec:prices}.
\item \texttt{code/calibrate.py}: builds peak-run energies from the archive of
\citet{sistig2025dataset}, written to \texttt{data/peak\_runs.json}.
\item \texttt{code/study.py}: the sweep, written to \texttt{results/summary.json}
and \texttt{results/study.json}; \texttt{code/figures.py}: all figures.
\item \texttt{make.sh} regenerates results, figures and the PDF.
\end{itemize}
\endgroup
The public-timetable certificates, their portable checkers and an index binding
sources to commands are under \texttt{research-20261029/} and
\texttt{paper/focused-theory-20261029/CERTIFICATE\_INDEX.md}.
